\documentclass[11pt]{article}
\usepackage[margin=1in]{geometry}
\usepackage{amsmath,amssymb,amsthm,mathtools}
\usepackage{graphicx}
\usepackage{microtype}
\usepackage[hidelinks]{hyperref}
\usepackage{xcolor}

\newtheorem{theorem}{Theorem}
\newtheorem{lemma}{Lemma}
\newtheorem{corollary}{Corollary}
\theoremstyle{definition}
\newtheorem{definition}{Definition}
\newtheorem{remark}{Remark}

\newcommand{\R}{\mathbb R}
\newcommand{\N}{\mathbb N}
\newcommand{\G}{\mathcal G}
\newcommand{\Phiop}{\Phi}
\newcommand{\norm}[1]{\left\lVert #1\right\rVert}

\title{Maximal unstable uniformity is independent of the complement\\
\large A candidate full negative resolution of Remark 39(b) in arXiv:2411.06209}
\author{Automated mathematics run \texttt{fa\_banach\_001}, lane 8}
\date{27 August 2026}

\begin{document}
\maketitle

\begin{center}
\fbox{\parbox{0.91\linewidth}{\textbf{Status.}
Candidate full negative solution, likely valid, pending expert review.  The
theorem below proves that the maximal unstable uniformity dimension is
independent of the complementary space in every dimension, thereby disproving
the existential conjecture in Remark 39(b) of the source paper.}}
\end{center}

\section{Source conjecture}

Czornik, Kitzing, and Siegmund study the one-sided system
\begin{equation}\label{eq:system}
 x(n+1)=A(n)x(n),\qquad n\in\N,
\end{equation}
where $A(n)\in\mathrm{GL}(\R^d)$ and both $A$ and $A^{-1}$ are bounded
\cite{CKS}.  On PDF and printed page 20, Remark 39(b) conjectures:
\begin{quote}
``There exists $A\colon\N\to\mathrm{GL}(\R^d)$ such that the maximal
uniformity dimension $u_2$ depends on the choice of $L_2$.''
\end{quote}
Here $L_2$ ranges over complements of the fixed stable space $L_1$.

\begin{figure}[ht]
 \centering
 \includegraphics[width=\linewidth]{figures/open_problem_crop.png}
 \caption{Theorem 38 and Remark 39 in the source paper, page 20.  Part (b) is
 the conjecture resolved below.}
\end{figure}

\section{Definitions and notation}

Write
\[
 \Phiop(0)=I,\qquad \Phiop(n)=A(n-1)\cdots A(0)\quad(n\ge1),
\]
so that the solution with initial value $x_0$ is $x(n,x_0)=\Phiop(n)x_0$.
A splitting is $\R^d=L_1\oplus L_2$.  For a subspace $L$, let
$\G_j(L)$ be its $j$-dimensional subspaces.

Following the source, the system has a \emph{dichotomy with uniformity
dimensions $(j_1,j_2)$} on $(L_1,L_2)$ if there is a single $\alpha>0$ such
that
\begin{align}
 \norm{\Phiop(n)u}&\le C_U e^{-\alpha(n-m)}\norm{\Phiop(m)u},
 &u&\in U,\quad U\in\G_{j_1}(L_1), \label{eq:stable}\\
 \norm{\Phiop(n)v}&\ge c_V e^{\alpha(n-m)}\norm{\Phiop(m)v},
 &v&\in V,\quad V\in\G_{j_2}(L_2), \label{eq:unstable}
\end{align}
for all $n\ge m\ge0$, where the positive constants may depend on the
subspace.  Refining a uniformity subspace preserves the estimate and its
exponent.  Thus a dichotomy is uniform on every one-dimensional subspace of
both splitting spaces \cite[Definition 15 and Theorem 7]{CKS}.

The pair $(u_1,u_2)$ is \emph{maximal} if it is an available uniformity pair
and every available pair $(\ell_1,\ell_2)$ satisfies
$\ell_i\le u_i$ \cite[Definition 35]{CKS}.

\section{Proof intuition}

An alternate complement is a graph over the original complement.  If
$w=v+Tv$ lies in that graph, then $v$ is in the old unstable space and $Tv$
is in the stable space.  The source proof tries to control $T(V)$ by a full
stable uniformity estimate, which leads to its extra dimension hypothesis.
That is more than is needed.

Indeed, finitely many stable basis-line estimates give the absolute bound
$\norm{\Phiop(n)s}\le K e^{-\alpha n}\norm{s}$ for every $s\in L_1$.
Meanwhile the old unstable estimate gives
$\norm{\Phiop(n)v}\ge c e^{\alpha n}\norm v$.  Hence the graph correction is
at most a constant times $e^{-2\alpha n}$ of the unstable component.  It is
therefore harmless after a finite time, and invertibility patches the finite
initial interval.

\section{Complement-transfer theorem}

\begin{lemma}[Finite-time patch]\label{lem:patch}
Let $U\subset\R^d$ be fixed.  Suppose that for some $C>0$, $\alpha>0$, and
$r_0\in\N$,
\[
 \norm{\Phiop(n)u}\ge C e^{\alpha(n-m)}\norm{\Phiop(m)u}
 \qquad(n\ge m\ge r_0,\ u\in U).
\]
Then the same estimate holds for all $n\ge m\ge0$ after replacing $C$ by a
possibly smaller positive constant.
\end{lemma}

\begin{proof}
Every $\Phiop(r)$ is invertible.  For each $m<r_0$, compactness of the unit
sphere of $U$ gives
\[
 q_m:=\inf_{0\ne u\in U}
 \frac{\norm{\Phiop(r_0)u}}{\norm{\Phiop(m)u}}>0.
\]
For $n\ge r_0>m$, apply the assumed estimate from $r_0$ to $n$ and insert
$q_m$; this gives the desired form with constant
$Cq_m e^{-\alpha(r_0-m)}$.  The remaining pairs $m\le n<r_0$ are finite in
number, and the corresponding norm ratios again have positive minima on the
unit sphere of $U$.  Taking the minimum of these finitely many constants
proves the claim.  This is also the mechanism of Lemma 36 in the source.
\end{proof}

\begin{theorem}[Uniformity dimensions transfer between all complements]
\label{thm:transfer}
Let $A\colon\N\to\mathrm{GL}(\R^d)$ satisfy the standing boundedness
assumptions of \eqref{eq:system}, and let
\[
 \R^d=L_1\oplus L_2=L_1\oplus L_2'.
\]
If the system has a dichotomy with uniformity dimensions $(j_1,j_2)$ on
$(L_1,L_2)$, then it has a dichotomy with the same uniformity dimensions on
$(L_1,L_2')$.  The same exponent $\alpha$ can be used.
\end{theorem}

\begin{proof}
The degenerate cases $L_1=\{0\}$ or $L_1=\R^d$ have a unique complement, so
assume $1\le k:=\dim L_1\le d-1$.  An admissible uniformity pair then has
$j_1,j_2\ge1$.

Let $e_1,\dots,e_k$ be an orthonormal basis of $L_1$.  By refinement of the
stable estimate, for each basis line there is $C_i>0$ such that, on setting
$m=0$ in \eqref{eq:stable},
\[
 \norm{\Phiop(n)e_i}\le C_i e^{-\alpha n}\qquad(n\ge0).
\]
Consequently, for $s=\sum_i a_i e_i$,
\begin{equation}\label{eq:absolute-stable}
 \norm{\Phiop(n)s}
 \le e^{-\alpha n}\sum_i |a_i|C_i
 \le K e^{-\alpha n}\norm{s},
 \qquad K:=\Big(\sum_i C_i^2\Big)^{1/2}.
\end{equation}
This is an absolute time-zero estimate on all of $L_1$; it does not assert
moving-time uniformity on $L_1$.

Fix an arbitrary $V'\in\G_{j_2}(L_2')$.  Let $P\colon\R^d\to L_2$ be the
projection along $L_1$.  Since $V'\cap L_1=\{0\}$, $P|_{V'}$ is injective,
and
\[
 V:=P[V']\in\G_{j_2}(L_2).
\]
Thus $V'$ is the graph of a linear map $T\colon V\to L_1$:
\[
 V'=\{v+Tv:v\in V\}.
\]
Put $M=\norm T$.  By \eqref{eq:unstable} on the fixed space $V$, there is
$c>0$ such that
\begin{equation}\label{eq:Vunstable}
 \norm{\Phiop(n)v}\ge c e^{\alpha(n-m)}\norm{\Phiop(m)v}
 \quad(n\ge m\ge0,\ v\in V).
\end{equation}
In particular, at $m=0$,
$\norm{\Phiop(r)v}\ge c e^{\alpha r}\norm v$.  Combining this with
\eqref{eq:absolute-stable}, for every $r\ge0$ and $v\in V$,
\begin{equation}\label{eq:ratio}
 \norm{\Phiop(r)Tv}
 \le KM e^{-\alpha r}\norm v
 \le D e^{-2\alpha r}\norm{\Phiop(r)v},
 \qquad D:=\frac{KM}{c}.
\end{equation}

Choose $r_0$ so that $D e^{-2\alpha r_0}\le\tfrac12$.  For
$w=v+Tv\in V'$ and $n\ge m\ge r_0$, equations \eqref{eq:Vunstable} and
\eqref{eq:ratio} give
\begin{align*}
 \norm{\Phiop(n)w}
 &\ge \norm{\Phiop(n)v}-\norm{\Phiop(n)Tv}
 \ge \tfrac12\norm{\Phiop(n)v},\\
 \norm{\Phiop(m)w}
 &\le \norm{\Phiop(m)v}+\norm{\Phiop(m)Tv}
 \le (1+D)\norm{\Phiop(m)v}.
\end{align*}
Therefore
\[
 \norm{\Phiop(n)w}
 \ge \frac{c}{2(1+D)}e^{\alpha(n-m)}\norm{\Phiop(m)w}
 \qquad(n\ge m\ge r_0,\ w\in V').
\]
Lemma~\ref{lem:patch} extends this to all $n\ge m\ge0$.  Hence the unstable
estimate holds uniformly on $V'$ with exponent $\alpha$.  Since $V'$ was an
arbitrary member of $\G_{j_2}(L_2')$, all required unstable estimates hold on
the new complement.  The stable estimates on $\G_{j_1}(L_1)$ are unchanged.
Thus $(j_1,j_2)$ transfers to $(L_1,L_2')$ with the same exponent.
\end{proof}

\begin{corollary}[Full negative answer to Remark 39(b)]
For one-sided time and fixed $L_1$, the maximal unstable uniformity dimension
$u_2$ is independent of the choice of complement $L_2$.  Hence the system
conjectured in Remark 39(b) of \cite{CKS} does not exist.
\end{corollary}

\begin{proof}
Let $u_2$ and $u_2'$ be the maximal unstable uniformity dimensions for two
complements $L_2$ and $L_2'$.  Theorem~\ref{thm:transfer} transfers the maximal
pair for $(L_1,L_2)$ to $(L_1,L_2')$, so maximality gives $u_2\le u_2'$.
Applying the same theorem in the reverse direction gives $u_2'\le u_2$.
Thus $u_2=u_2'$.
\end{proof}

\section{Why the source hypothesis is unnecessary}

Theorem 37 of \cite{CKS} proves $u_2\le u_2'$ under the additional condition
that the stable projection of every target $u_2$-plane has dimension at most
$u_1$.  That condition supplies a moving-time stable estimate on the entire
projected plane.  The proof above replaces it with
\eqref{eq:absolute-stable}, obtained from finitely many basis lines and used
only at absolute time zero.  This is sufficient because the stable and
unstable absolute-time bounds separate at the rate $e^{-2\alpha r}$.  No
dimension restriction on the stable projection remains.

\section{Verification, novelty, and limitations}

The proof uses only finite-dimensional linear algebra, the definitions and
refinement observation in the source, and the explicitly proved finite-time
patch.  The constants $c$, $D$, and $r_0$ may depend on the chosen target
plane $V'$, exactly as the source definition permits; the exponent $\alpha$
does not.  Degenerate splittings have a unique complement.

A bounded novelty check on 27 August 2026 searched the run's lightweight
indexes and live arXiv/web results by the exact arXiv id, paper title, exact
conjecture language, authors, and the phrases ``maximal uniformity dimension'',
``dependence on splitting'', and ``choice of $L_2$''.  It found the source
paper and a repository mirror, but no later claimed resolution.  This is not
an exhaustive novelty guarantee.

There is no computational component.  Human review should focus on the common
exponent under refinement, the graph parametrization of an arbitrary target
$j_2$-plane, and the fact that only the absolute stable estimate
\eqref{eq:absolute-stable} is used.

\enlargethispage{5\baselineskip}
\begin{thebibliography}{9}
\bibitem{CKS}
A.~Czornik, K.~Kitzing, and S.~Siegmund,
\emph{Dichotomies uniform on subspaces and formulas for dichotomy spectra},
arXiv:2411.06209, latest source/PDF inspected 27 August 2026.
\end{thebibliography}

\end{document}
